\chapter{Fundamentos Teoricos y Estado del Arte}

\section{2.1 Entrenamiento de redes profundas y grafos computacionales}
Este apartado establece la semantica de ejecucion por grafo y su relacion con dependencias de datos. La descripcion sirve de base para representar costos por nodo y por arista.

\section{2.2 Arquitecturas heterogeneas CPU-GPU y jerarquias de memoria}
Se discute la complementariedad entre CPU y GPU en terminos de computo, memoria y latencia. La jerarquia de memoria se incorpora como restriccion central del problema de particion.

\section{2.3 Transferencia de datos, solapamiento y comunicacion}
Se analizan costos de transferencia inter-dispositivo y su posible solapamiento con computo. Esta seccion fundamenta el termino de costo de corte en la formulacion de optimizacion.

\section{2.4 Estrategias de gestion de memoria}
Se revisan enfoques de precision, recomputo y persistencia de activaciones que condicionan viabilidad de entrenamiento bajo presupuestos acotados.

\section{2.5 Optimizacion declarativa: ILP en sistemas computacionales}
Se presenta el uso de ILP como marco de decision con restricciones explicitas y objetivo multicriterio, destacando trazabilidad de solucion y capacidad de auditoria.

\section{2.6 Estado del arte y vacios de investigacion}
Se sintetizan contribuciones previas y se delimita el vacio abordado por esta tesis: cierre metodologico entre modelado, simulacion y validacion operacional reproducible.
